\chapter{Piedra, papel o tijera}
\label{chap:piedrapapeltijera}
Este es el juego de piedra, papel o tijera de la aplicación.
El jugador elige una jugada, la aplicación elige otra al azar y se muestra quién gana.
Un marcador cuenta las victorias, las derrotas y los empates.

Se sigue la misma paleta de colores que en el resto de la aplicación.
Las jugadas, el marcador y el resultado se muestran en fichas con borde petróleo, como las del menú.

\begin{figure}[H]
  \centering
  \includegraphics[width=0.4\textwidth]{piedrapapeltijera}
\end{figure}

\section{Recursos visuales}
Los iconos de las tres jugadas y el signo de interrogación inicial son \texttt{VectorDrawable} en la carpeta \texttt{drawable}.
Se llaman \texttt{ic\_piedra}, \texttt{ic\_papel}, \texttt{ic\_tijera} e \texttt{ic\_pregunta}.
Al ser vectores se ven nítidos a cualquier tamaño y usan los colores de la paleta de la aplicación.
El interrogante solo ocupa el sitio de las dos jugadas hasta la primera partida.

\section{Representación de las jugadas}
Cada jugada es un número: $0$ es piedra, $1$ es papel y $2$ es tijera.
Una lista con los tres iconos, en ese mismo orden, permite mostrar la jugada con su número.

\begin{lstlisting}
private val iconos = listOf(R.drawable.ic_piedra, R.drawable.ic_papel, R.drawable.ic_tijera)
private val nombres = listOf("piedra", "papel", "tijera")
\end{lstlisting}

\section{Elegir al ganador}
La aplicación elige su jugada con un número aleatorio entre $0$ y $2$.
Después se compara con la del jugador.

\begin{lstlisting}
val app = (0..2).random()

if (jugador == app) {
    empates++
} else if ((jugador == 0 && app == 2) || (jugador == 1 && app == 0) || (jugador == 2 && app == 1)) {
    victorias++
} else {
    derrotas++
}
\end{lstlisting}

Hay empate si las dos jugadas son iguales.
El jugador gana si la piedra vence a la tijera, el papel a la piedra o la tijera al papel.
En cualquier otro caso, gana la aplicación.

\section{Resultado y marcador}
Después de cada ronda se ponen los dos iconos con \texttt{setImageResource}.
El texto del resultado dice qué jugada eligió la aplicación.
El color del banner cambia según el resultado: verde si gana el jugador, naranja si pierde y mostaza si empata.

El marcador son tres contadores que se actualizan al terminar cada ronda.
El botón Reiniciar marcador los pone a cero y devuelve los dos iconos al interrogante.